% Thesis-ready fragment for inclusion in the Outlook chapter.

\subsection{An on-shell formulation of the Michel parameters}

A natural continuation of the present analysis would be to formulate the
new-physics perturbations of the Michel parameters directly in an on-shell
amplitude basis. Instead of beginning with a Lagrangian operator basis or the
traditional SPVAT couplings, one would construct the most general local
four-point amplitude for muon decay,
\begin{equation}
    \mathcal A\!\left(\mu\to e\nu_p\bar\nu_r\right)
    = \mathcal A_{\mathrm{SM}}
    + \sum_I a_I\,\mathcal B_I\,.
    \label{eq:outlook-onshell-amplitude}
\end{equation}
Here, the $\mathcal B_I$ denote independent spinor-helicity structures that
obey Lorentz invariance, little-group covariance, locality and the required
fermionic exchange symmetries. The muon and, for a complete analysis of the
Michel parameters, the electron would have to be treated as massive external
states, whereas the neutrino masses could be neglected. The coefficients
$a_I$ would constitute an amplitude-level counterpart of the LEFT Wilson
coefficients and could be related to them by an explicit matching calculation.

The Michel parameters would then arise from bilinears of these amplitude
coefficients after forming the polarized squared amplitude and integrating
over the unobserved neutrino phase space. Schematically, each parameter could
be represented as
\begin{equation}
    X
    = \frac{\boldsymbol a^{\dagger} M_X\boldsymbol a}
           {\boldsymbol a^{\dagger} M_N\boldsymbol a}\,,
    \qquad
    X\in\{\rho,\eta,\xi,\delta,\xi',\ldots\}\,,
    \label{eq:outlook-michel-bilinear}
\end{equation}
where $M_X$ encodes the relevant polarized phase-space integral and $M_N$
determines the decay normalization. Expanding
$\boldsymbol a=\boldsymbol a_{\mathrm{SM}}+\delta\boldsymbol a$ would directly
separate the interference terms, which are linear in the new amplitude, from
the contributions that are quadratic in it. In this way, the perturbative
expressions derived in the LEFT could be reproduced without introducing
off-shell operator redundancies.

Such a formulation would make several structural properties of the present
results particularly transparent. A new amplitude can interfere with the
Standard Model only if it produces the same neutrino-flavor and helicity state;
amplitudes with orthogonal invisible final states therefore first contribute
quadratically. Interference between different charged-lepton chiralities is
likewise absent in the massless limit or suppressed by the electron mass. The
latter must consequently be retained in order to recover the Michel parameter
$\eta$ and the transverse-polarization observables. For lepton-number-violating
channels, the antisymmetrization of identical neutrinos could be imposed
directly at amplitude level. This would provide a particularly transparent
interpretation of the symmetric and mixed flavor components of
$C_{\bar e LLLH}$. Moreover, operator expressions related by Fierz identities
would correspond to the same physical on-shell structure rather than appearing
as distinct intermediate descriptions.

The on-shell formulation would also establish a direct connection to the UV
dictionary employed in this thesis. A local four-point interaction observed at
low energies can be opened into possible ultraviolet origins by determining
the massive poles through which the corresponding amplitude can factorize. In
the vicinity of a heavy propagator, one schematically has
\begin{equation}
    \mathcal A_4^{\mathrm{UV}}
    \supset
    \frac{\mathcal A_3^{(L)}\mathcal A_3^{(R)}}{p^2-M^2}
    \xrightarrow{\,p^2\ll M^2\,}
    -\frac{\mathcal A_3^{(L)}\mathcal A_3^{(R)}}{M^2}
    +\mathcal O\!\left(\frac{p^2}{M^4}\right)\,.
    \label{eq:outlook-onshell-factorization}
\end{equation}
The low-energy contact amplitude is thus the first term in the expansion of a
factorization channel containing a dynamical heavy field. This would express
the full bottom-up trajectory---from Michel-parameter measurements, through
LEFT and SMEFT amplitudes, to possible ultraviolet particles and their
three-point couplings---within a common on-shell language.

This reformulation would not, however, eliminate the phenomenological part of
the calculation. The construction of polarized spin-density matrices, the
integration over the invisible phase space and the treatment of the nonzero
electron mass would still be required. A precision comparison with experiment
would additionally have to incorporate QED radiative corrections and the
experimental definition of the measured decay distributions. Developing such
a framework therefore represents a substantial project rather than merely a
change of notation. Nevertheless, it could reveal the helicity and
factorization structure underlying the Michel parameters more directly and
provide a conceptually unified continuation of the bottom-up analysis pursued
in this thesis.

The spinor-helicity formalism also points beyond the current Lagrangian formulation of particle physics. This fits very well with the overall directive of this work, namely abandoning the idea of a single, complete theory that transcends the observables but rather display the principles that connect the multiple description of physical reality.